\exercisesection{}

¶ 1) Let G be a locally compact group, β a left Haar measure on G, A a β-measurable subset of G, and ν a nonzero positive measure on G. Assume that sA is locally ν-negligible for every \(s \in G\) . Show that A is locally β-negligible. (Reduce to the case that A is relatively compact and ν is bounded. Prove that ν and \(\varphi_{A} \cdot \beta\) are convolvable and that \(\nu *^{\beta} \varphi_{A} = 0\) , whence \(0 = \|\nu * \varphi_{A} \cdot \beta\| = \|\nu\| \cdot \|\varphi_{A} \cdot \beta\|\) .) Show, on admitting the continuum hypothesis, that this result may fail to hold if A is not assumed to be β-measurable. (Take \(G = R^{2}\) , take for ν the Haar measure on the subgroup \(R \times \{0\}\) , and apply Exer. 7 c) of Ch. V, §8.)

\begin{enumerate}
\def\labelenumi{\arabic{enumi})}
\setcounter{enumi}{1}
\tightlist
\item
  Let H be the additive group R equipped with the discrete topology. Let G be the locally compact group \(R \times R \times H\) . Let \(\alpha\) be a Haar measure on G, \(\beta\) a Haar measure on \(R \times H\) , and \(\mu = \varepsilon_{0} \otimes \beta \in \mathcal{M}(G)\) . Construct a function \(f \geqslant 0\) on G, locally \(\alpha\) -negligible (hence such that \(\mu\) and f are convolvable) but such that no translate of f is \(\mu\) -measurable. (Imitate the construction of Ch. V, §3, Exer. 4.)
\end{enumerate}

¶ 3) Let G be a locally compact group.

\begin{enumerate}
\def\labelenumi{\alph{enumi})}
\tightlist
\item
  Let \(\mu\) be a nonzero bounded positive measure on \(\mathbf{G}\) such that \(\mu * \mu = \mu\) . Show that the support \(\mathbf{S}\) of \(\mu\) is compact. (Let \(f \in \mathcal{K}_{+}(\mathbf{G})\) with \(f \neq 0\) ; choose a left Haar measure on \(\mathbf{G}\) , with respect to which convolution products will be taken; we have \(\mu * f \in \overline{\mathcal{K}(\mathbf{G})}\) ; let \(x \in \mathbf{S}\) be such that \((\mu * f)(x) = \sup_{y \in \mathbf{S}} (\mu * f)(y)\) ; show that
\end{enumerate}

\((\mu * f)(y) = (\mu * f)(x)\) for every \(y \in S\) .

\begin{enumerate}
\def\labelenumi{\alph{enumi})}
\setcounter{enumi}{1}
\tightlist
\item
  Show that S is a compact subgroup of G and that \(\mu\) is the normalized Haar measure of S. (Make use of a), Exer. 21 of GT, III, §4, and Exer. 5 of §3.)
\end{enumerate}

\begin{enumerate}
\def\labelenumi{\arabic{enumi})}
\setcounter{enumi}{3}
\tightlist
\item
  Let \(G\) be a locally compact group. For every \(t \in \mathbf{R}_+^*\) , let \(\mu_t\) be a nonzero bounded positive measure on \(G\) . Assume that the mapping \(t \mapsto \mu_t\) is continuous for the topology \(\sigma(\mathcal{M}^1(G), \overline{\mathcal{K}(G)})\) , and that \(\mu_{s+t} = \mu_s * \mu_t\) ( \(s, t\) in \(\mathbf{R}_+^*\) ).
\end{enumerate}

\begin{enumerate}
\def\labelenumi{\alph{enumi})}
\item
  Show that there exists a number \(c \in \mathbf{R}\) such that \(\| \mu_t \| = \exp(ct)\) . (Observe that \(t \mapsto \| \mu_t \|\) is lower semi-continuous and that \(\| \mu_{s+t} \| = \| \mu_s \| \cdot \| \mu_t \|\) . Make use of Prop. 18.)
\item
  Suppose that \(c = 0\) . Show that \(\mu_t\) converges weakly as \(t \to 0\) to the normalized Haar measure of a compact subgroup of \(G\) . (Make use of the weak compactness of the unit ball of \(\mathcal{M}^1(G)\) , and show that for every weak cluster point \(\mu\) of \(t \mapsto \mu_t\) with respect to the filter of neighborhoods of 0 in \(R_+^*\) , one has \(\mu_t * \mu = \mu_t\) for every \(t\) , then \(\mu * \mu = \mu\) ; next, apply Exer. 3.)
\end{enumerate}

¶ 5) Let G be a locally compact group countable at infinity, β a left Haar measure on G, \(a \in R_{+}^{*}\) , and \((\nu_{u})_{0 \leqslant u \leqslant a}\) a family of positive measures on G satisfying the following conditions:

\begin{enumerate}
\def\labelenumi{(\roman{enumi})}
\item
  \(\nu_0 = \varepsilon_e\) ; \(\nu_{u + v} = \nu_u * \nu_v\) if \(u + v \leqslant a\) ; \(\nu_u = \stackrel{\vee}{\nu}_u\) ;
\item
  for \(0 < u \leqslant a\) , \(\nu_{u} = f_{u} \cdot \beta\) with \(f_{u} \geqslant 0\) lower semi-continuous;
\item
  \(\nu_{u}\) is a vaguely continuous function of \(u\) .
\end{enumerate}

\begin{enumerate}
\def\labelenumi{\alph{enumi})}
\tightlist
\item
  Let \(f \in \mathcal{K}_{+}(\mathrm{G})\) . Show that, for \(0 < u \leqslant a\) ,
\end{enumerate}

\[
\int \left(f _ {u / 2} * f\right) (z) ^ {2} d \beta (z) = \int f (x) \left(f _ {u} * f\right) (x) d \beta (x).\tag{1}
\]

\begin{enumerate}
\def\labelenumi{\alph{enumi})}
\setcounter{enumi}{1}
\item
  Let \(f \in \mathcal{K}(\mathbf{G})\) . Show that \(f_{u/2} * f\) has \(\beta\) -integrable square for \(0 < u \leqslant a\) , and that (1) again holds.
\item
  Let \(f \in \mathcal{H}(\mathrm{G})\) be such that \(\nu_{a} * f = 0\) . Show that \(f = 0\) . (One has \(\nu_{a/2^n} * f = 0\) for every integer \(n > 0\) by \(b\) ), therefore \(f = 0\) by §3, No.~3, Remark 1).
\item
  Let \(\nu \in \mathcal{C}'(\mathrm{G})\) be such that \(\nu_{a} * \nu = 0\) . Show that \(\nu = 0\) . (Regularize \(\nu\) by functions in \(\mathcal{H}(\mathrm{G})\) and apply \(c\) ).)
\end{enumerate}

\begin{enumerate}
\def\labelenumi{\arabic{enumi})}
\setcounter{enumi}{5}
\item
  Let G be a locally compact group. Show that if \(\mathcal{K}(\mathrm{G})\) is commutative for convolution, then G is abelian. (Show by regularization that \(\mathcal{C}'(\mathrm{G})\) is commutative, and apply this to the measures \(\varepsilon_{s}\) , where \(s \in G\) .)
\item
  Let \(G\) be a locally compact group and \(\beta\) a left Haar measure on \(G\) . Show that the algebra \(L^1(G, \beta)\) has a unity element if and only if \(G\) is discrete. (Suppose \(G\) non-discrete, and let \(f_0 \in L^1(G, \beta)\) . There exists a compact neighborhood \(V\) of \(e\) such that
\end{enumerate}

\[
\int_ {V} | f _ {0} (x) | d \beta (x) <   1.
\]

Let \(\mathbf{U}\) be a symmetric compact neighborhood of \(e\) such that \(\mathbf{U}^2 \subset \mathbf{V}\) . Then, for almost every \(x \in \mathbf{U}\) ,

\[
\left| \left(\varphi_ {\mathrm{U}} * f _ {0}\right) (x) \right| = \int_ {\mathrm{U}} \left| f _ {0} \left(y ^ {- 1} x\right) \right| d \beta (y) \leqslant \int_ {\mathrm{V}} \left| f _ {0} (x) \right| d \beta (x) <   1,
\]

thus \(f_0\) is not a unity element for \(L^1(G, \beta)\) .

\begin{enumerate}
\def\labelenumi{\arabic{enumi})}
\setcounter{enumi}{7}
\tightlist
\item
  Let G be a locally compact group countable at infinity, operating continuously on the left in a Polish locally compact space T. Let \(\nu\) be a positive measure on T that is quasi-invariant under G. Let R be a \(\nu\) -measurable equivalence relation on T, compatible with G. There then exists (Ch. VI, §3, No.~4, Prop. 2) a Polish locally compact space B, and a \(\nu\) -measurable mapping p of T into B, such that \(R\{x,y\}\) is equivalent to \(p(x)=p(y)\) . Let \(\nu'\) be a pseudo-image measure of \(\nu\) under p, and let \(b\mapsto\lambda_{b}\) ( \(b\in B\) ) be a disintegration of \(\nu\) by R. Show that the \(\lambda_{b}\) are, for almost every \(b\in B\) , quasi-invariant under G.
\end{enumerate}

(Let \(\chi\) be a function on \(\mathbf{G} \times \mathbf{T}\) satisfying the conditions of Exer. 13 of Ch. VII, §1. Show that for every \(s \in \mathbf{G}\) , there exists a \(\nu'\) -negligible subset \(\mathbf{N}(s)\) of B such that \(\chi(s^{-1}, \cdot)\) is locally \(\lambda_b\) -integrable for \(b \notin \mathbf{N}(s)\) ; for this, observe that for every \(\psi \in \mathcal{K}(\mathbf{T})\) , the function \(x \mapsto \psi(x)\chi(s^{-1}, x)\) is \(\lambda_b\) -integrable except for \(b\) belonging to a \(\nu'\) -negligible set \(\mathbf{N}(s, \psi)\) , and make use of Lemma 1 of Ch. VI, §3, No.~1. Set \(\lambda_{b,s}' = 0\) if \(b \in \mathbf{N}(s)\) , and \(\lambda_{b,s}' = \chi(s^{-1}, \cdot) \cdot \lambda_b\) if \(b \notin \mathbf{N}(s)\) . Show that the mapping \(b \mapsto \lambda_{b,s}'\) is \(\nu'\) -adequate (use Lemma 3 of Ch. VI, §3, No.~1) and that \(\gamma(s)\beta = \int \lambda_{b,s}' d\nu'(b)\) . Show that on the other hand \(\gamma(s)\beta = \int \gamma(s)\lambda_b d\nu'(b)\) and deduce from this that, for every \(s \in \mathbf{G}\) , one has \(\gamma(s)\lambda_b = \chi(s^{-1}, \cdot) \cdot \lambda_b\) for almost every \(b\) , therefore for almost every \(b\) one has \(\gamma(s)\lambda_b = \chi(s^{-1}, \cdot) \cdot \lambda_b\) for almost every \(s\) . Then use Cor. 2 of Prop. 17 of §4 to infer that, for almost every \(b\) , \(\gamma(s)\lambda_b\) is equivalent to \(\lambda_b\) for all \(s \in \mathbf{G}\) .)

Show that if \(\nu\) is relatively invariant under G with multiplier \(\chi\) , then the \(\lambda_{b}\) are, for almost every b, relatively invariant with multiplier \(\chi\) .

\begin{enumerate}
\def\labelenumi{\arabic{enumi})}
\setcounter{enumi}{8}
\item
  Let \(G\) be a locally compact group, \(\beta\) a left Haar measure on \(G\) , \(A\) and \(B\) two \(\beta\) -integrable sets such that \(\beta(A) \leqslant \beta(B)\) and \(\beta^{*}(A^{-1}) < +\infty\) . Show that there exist
\item
  disjoint sets \(N, K_{1}, K_{2}, \ldots\) covering \(A\) , with \(N \beta\) -negligible and the \(K_{n}\) compact;
\item
  disjoint sets \(\mathbf{N}^{\prime}\) , \(\mathbf{K}_1^{\prime}\) , \(\mathbf{K}_2^{\prime}\) , \ldots{} covering B, with the \(\mathbf{K}_n^{\prime}\) compact;\\
\item
  elements \(s_n \in G\) such that \(\mathbf{K}_n^{\prime} = s_n\mathbf{K}_n\) .
\end{enumerate}

(Using the fact that \(\beta (x\mathbf{A}\cap \mathbf{B})\) depends continuously on \(x\) and that

\[
\int \beta (x \mathrm{A} \cap \mathrm{B}) d \beta (x) = \beta (\mathrm{A} ^ {- 1}) \beta (\mathrm{B}),
\]

show that if \(\beta (\mathbf{A})\neq 0\) then there exists an \(x\in \mathbf{G}\) such that \(\beta (x\mathbf{A}\cap \mathbf{B})\neq 0\) .)

\begin{enumerate}
\def\labelenumi{\arabic{enumi})}
\setcounter{enumi}{9}
\tightlist
\item
  Let \(\mathbf{G}\) be a locally compact group, \(\beta\) a left Haar measure on \(\mathbf{G}\) . For any two \(\beta\) -integrable sets \(\mathbf{A}, \mathbf{B}\) set \(\rho(\mathbf{A}, \mathbf{B}) = \beta(\mathbf{A} \cup \mathbf{B}) - \beta(\mathbf{A} \cap \mathbf{B})\) .
\end{enumerate}

\begin{enumerate}
\def\labelenumi{\alph{enumi})}
\item
  Let A be a \(\beta\) -integrable set. Show that \(x \mapsto \rho(xA, A)\) is a continuous function.
\item
  Let U be a neighborhood of e. Show that there exist a compact set A and a number \(\varepsilon > 0\) such that \(\rho(x\mathbf{A}, \mathbf{A}) < \varepsilon\) implies \(x \in U\) . (Take for A a neighborhood of e such that \(A \cdot A^{-1} \subset U\) .)
\item
  For a subset C of G to be relatively compact, it is necessary and sufficient that there exist a \(\beta\) -integrable set A and a number \(a (0 < a < 2\beta(A))\) such that \(x \in C\) implies \(\rho(xA, A) \leqslant a\) .
\end{enumerate}

\begin{enumerate}
\def\labelenumi{\arabic{enumi})}
\setcounter{enumi}{10}
\tightlist
\item
  \begin{enumerate}
  \def\labelenumii{\alph{enumii})}
  \tightlist
  \item
    Let G be a locally compact group generated by a compact neighborhood of e. Let \(\varphi\) be a non-surjective continuous endomorphism of G belonging to the closure, for the topology of compact convergence, of the group G of (bicontinuous) automorphisms of G. Then \(\lim_{\psi\in\mathcal{G},\psi\to\varphi}\mod\psi=0\) . (Let K be a compact neighborhood of e that generates G. Let \(\mu\) be a left Haar measure on G. If \(\mu(\varphi(K)) > 0\) , then \(\varphi(K) \cdot \varphi(K)^{-1}\) is a neighborhood of \(e\) in G, therefore \(\varphi(G)\) is an open subgroup of G; for \(\psi \in \mathcal{G}\) sufficiently near \(\varphi\) , one has \(\psi(K) \subset \varphi(G)\) , therefore \(\psi(G) \neq G\) , which is absurd. Thus \(\mu(\varphi(K)) = 0\) . As \(\psi \in \mathcal{G}\) tends to \(\varphi\) , \(\mu(\psi(K))\) tends to 0.)
  \end{enumerate}
\end{enumerate}

\begin{enumerate}
\def\labelenumi{\alph{enumi})}
\setcounter{enumi}{1}
\tightlist
\item
  Let G be a free abelian group, a direct sum \(G_{1} \oplus G_{2} \oplus \cdots\) , where each \(G_{i}\) is isomorphic to Z. Consider G as being discrete. Let \(\varphi\) be the non-surjective endomorphism \((x_{1}, x_{2}, \ldots) \mapsto (0, x_{1}, x_{2}, \ldots)\) of G. Then \(\varphi\) is the limit, for the topology of compact convergence, of automorphisms of G, and every automorphism of G has modulus 1.
\end{enumerate}

\begin{enumerate}
\def\labelenumi{\arabic{enumi})}
\setcounter{enumi}{11}
\tightlist
\item
  For \(t > 0\) and \(x \in \mathbf{R}\) , let \(F_t(x) = te^{-\pi t^2 x^2}\) . Let \(f \in \overline{\mathcal{H}(\mathbf{R})}\) . Show that \(f * F_t\) tends to \(f\) uniformly as \(t\) tends to \(+\infty\) . (Show that the measures \(F_t(x) dx\) satisfy the conditions of §2, No.~7, Lemma 4, with \(a = 0\) .)
\end{enumerate}

¶ 13) Let G be a locally compact group operating continuously on the left in a Polish locally compact space T. Let β be a left Haar measure on G, ν a quasi-invariant positive measure on T, and χ(s,x) a function \textgreater0 on G × T satisfying the conditions of Exer. 13 of Ch. VII, §1.

For \(f \in \mathrm{L}^p(\mathrm{T}, \nu)\) ( \(1 \leqslant p < +\infty\) ), set

\[
\left(\gamma_ {\chi , p} (s) f\right) (x) = \chi (s ^ {- 1}, x) ^ {1 / p} f (s ^ {- 1} x).
\]

\begin{enumerate}
\def\labelenumi{\alph{enumi})}
\item
  Show that for every \(s \in G\) , \(\pmb{\gamma}_{\chi,p}(s)\) is an isometric endomorphism of \(L^p(T, \nu)\) . Show that the mapping \(s \mapsto \pmb{\gamma}_{\chi,p}(s)\) is a linear representation of \(G\) in \(L^p(T, \nu)\) (argue as in §2, No.~5).
\item
  Let \(f \in \mathcal{L}^p(\mathbf{T}, \nu)\) , and let \(h \in \mathcal{K}(\mathbf{G})\) . Show that the function
\end{enumerate}

\[
(x, s) \mapsto f (s ^ {- 1} x) h (s) \chi (s ^ {- 1}, x) ^ {1 / p}
\]

is \(p\) -th power integrable for \(\beta \otimes \nu\) (begin with the case \(p = 1\) , and use Lemma 1 of §4, No.~1). Show that if \(q\) is the exponent conjugate to \(p\) , and if \(g \in \mathcal{L}^q(\mathrm{T}, \nu)\) , then \(f(s^{-1}x)h(s)g(x)\chi(s^{-1},x)^{1/p}\) is integrable for \(\beta\otimes\nu\) (write \(h\) in the form \(h_1h_2\) with \(h_1,h_2\) in \(\mathcal{K}(\mathrm{G})\) ). Then deduce from the Lebesgue-Fubini theorem that, for \(f\in\mathrm{L}^{p}(\mathrm{T},\nu)\) and \(g\in\mathrm{L}^{q}(\mathrm{T},\nu)\) , the function

\[
s \mapsto \int g (x) (\gamma_ {\chi , p} (s) f) (x) d \nu (x)
\]

is \(\beta\) -measurable.

\begin{enumerate}
\def\labelenumi{\alph{enumi})}
\setcounter{enumi}{2}
\item
  Show, using Cor. 2 of Prop. 18 of §4, No.~6 and Lemma 1 of Ch. VI, §3, No.~1, that the representation \(s \mapsto \gamma_{\chi,p}(s)\) of G in \(L^p(T, \nu)\) is continuous for \(1 \leqslant p < +\infty\) .
\item
  Assume in addition that, for every \(s \in G\) , the function \(x \mapsto \chi(s^{-1}, x)\) is bounded. For \(f \in L^{p}(T, \nu)\) , set
\end{enumerate}

\[
\left(\gamma_ {\chi} (s) f\right) (x) = \chi (s ^ {- 1}, x) f (s ^ {- 1} x).
\]

Show that \(s \mapsto \gamma_{\chi}(s)\) is a continuous representation of \(G\) in \(L^p(T, \nu)\) for \(1 \leqslant p < +\infty\) . (One shows, as in the proof of Prop. 9 of §2, No.~5, that \(s \mapsto \gamma_{\chi}(s)\) is a representation of \(G\) by endomorphisms of \(L^p(T, \nu)\) . One then observes that if \(f \in L^p(T, \nu)\) and \(h \in \mathcal{K}(G)\) , Lemma 1 of §4, No.~1 shows that \(h(s)f(s^{-1}x)\chi(s^{-1}, x)\) is \(p\) -th power integrable for \(\beta \otimes \nu\) . The proof is concluded as in c).)

\begin{enumerate}
\def\labelenumi{\arabic{enumi})}
\setcounter{enumi}{13}
\tightlist
\item
  \begin{enumerate}
  \def\labelenumii{\alph{enumii})}
  \tightlist
  \item
    Let \(G\) be a locally compact group, \(f\) a lower semi-continuous positive function on \(G\) , \(\mu\) a positive measure on \(G\) . Show that the function
  \end{enumerate}
\end{enumerate}

\[
x \mapsto \int^ {*} f (s ^ {- 1} x) d \mu (s) = g (x)
\]

is lower semi-continuous on G (cf.~Ch. IV, §1, No.~1, Th. 1); for \(\mu\) and \(f\) to be convolvable, it is necessary and sufficient that \(g\) be integrable for a left Haar measure on G.

\begin{enumerate}
\def\labelenumi{\alph{enumi})}
\setcounter{enumi}{1}
\tightlist
\item
  On the group \(\mathbf{G} = \mathbf{R} \times \mathbf{R}\) , let \(\mu\) be the measure \(\varepsilon_0 \otimes \lambda\) , where \(\lambda\) is Lebesgue measure on \(\mathbf{R}\) ; let \(f(x, y) = (1 - |xy - 2|)^+\) ; show that the function
\end{enumerate}

\[
g (x, y) = \int f (x - s, y - t) d \mu (s, t)
\]

is everywhere finite on G, but is not continuous and is not integrable for Lebesgue measure on G.

¶ 15) Let G be a locally compact group, β a left Haar measure on G; by an abuse of language, in what follows one identifies β-integrable numerical functions with their classes in L¹(G,β); same abuse for the Lᵖ(G,β).

\begin{enumerate}
\def\labelenumi{\alph{enumi})}
\item
  Let \(A\) be a continuous endomorphism of the Banach space \(\mathrm{L}^1(\mathrm{G},\beta)\) such that, for every \(s \in \mathrm{G}\) , one has \(A(f*\varepsilon_s) = A(f)*\varepsilon_s\) for all \(f \in \mathrm{L}^1(\mathrm{G},\beta)\) . Show that for every function \(g \in \mathcal{K}(\mathrm{G})\) , one then has \(A(f*g) = A(f)*g\) (observe that \(s \mapsto g(s)A(f*\varepsilon_s)\) is a \(\beta\) -integrable mapping of \(\mathrm{G}\) into \(\mathrm{L}^1(\mathrm{G},\beta)\) ); converse. From this, deduce that there exists one and only one bounded measure \(\mu\) on \(\mathrm{G}\) such that \(A(f) = \mu *f\) for all \(f \in \mathrm{L}^1(\mathrm{G},\beta)\) , and that \(\| A\| = \| \mu \|\) . (With the notations of Prop. 19 of No.~7, consider the limit of \(A(f_{\mathrm{V}}*g)\) with respect to an ultrafilter finer than the section filter of \(\mathfrak{B}\) , making use of the compactness of the unit ball of \(\mathcal{M}^1(\mathrm{G})\) for the weak topology \(\sigma (\mathcal{M}^1(\mathrm{G}),\overline{\mathcal{K}(\mathrm{G})})\) .)
\item
  Let \(\mu\) be a bounded measure on G; show directly that the norm of the continuous endomorphism \(\pmb{\gamma}(\mu): f \mapsto \mu * f\) of \(\mathrm{L}^{\infty}(\mathrm{G},\beta)\) is equal to \(\| \mu \|\) . (Reduce to the case that \(\mu\) has compact support and has a continuous density with respect to \(|\mu|\) .) From this, deduce anew that the continuous endomorphism \(\pmb{\gamma}(\mu): f \mapsto \mu * f\) of \(\mathrm{L}^{1}(\mathrm{G},\beta)\) has norm equal to \(\| \mu \|\) .
\item
  Assume that G is compact and \(\mu\) is positive. Show that for \(1 < p < +\infty\) , the norm of the continuous endomorphism \(\gamma(\mu): f \mapsto \mu * f\) of \(\mathrm{L}^{p}(\mathrm{G}, \beta)\) is equal to \(\|\mu\|\) .
\item
  Take for G the cyclic group of order 3. Give an example of a measure \(\mu\) on G such that the norm of the endomorphism \(\pmb{\gamma}(\mu)\) of \(\mathrm{L}^p (\mathrm{G},\beta)\) is strictly less than \(\| \mu \|\) for \(1 < p < +\infty\) .
\end{enumerate}

¶ 16) Notations and conventions are those of Exer. 15.

\begin{enumerate}
\def\labelenumi{\alph{enumi})}
\tightlist
\item
  Show that, for a bounded measure \(\mu\) on \(\mathbf{G}\) to be such that \(\| \mu * f \|_1 = \| f \|_1\) for all \(f \in \mathrm{L}^1(\mathbf{G}, \beta)\) , it is necessary and sufficient that \(\mu\) be a point measure of norm 1. (Using the fact that the endomorphism \(\gamma(|\mu|)\) of \(\mathrm{L}^1(\mathbf{G}, \beta)\) has norm \(\| \mu \|\) (Exer. 15), show that one necessarily has, for every function \(f \in \mathcal{K}(\mathbf{G})\) ,
\end{enumerate}

\[
\left| \int f d \mu \right| = \int | f | d | \mu |;
\]

from this, first deduce that \(\mu = c|\mu|\) , where \(c\) is a constant of absolute value 1, then that \(\mu\) is a point measure.)

\begin{enumerate}
\def\labelenumi{\alph{enumi})}
\setcounter{enumi}{1}
\tightlist
\item
  Take for G the cyclic group of order 3. Give an example of a measure \(\mu\) on G, not a point measure, such that \(\| \mu * f \|_2 = \| f \|_2\) for every numerical function \(f\) defined on G.
\end{enumerate}

\begin{enumerate}
\def\labelenumi{\arabic{enumi})}
\setcounter{enumi}{16}
\tightlist
\item
  Notations and conventions are those of Exer. 15.
\end{enumerate}

\begin{enumerate}
\def\labelenumi{\alph{enumi})}
\item
  Let \(\mu\) be a bounded measure on G; for the endomorphism \(\pmb{\gamma}(\mu): f \mapsto \mu * f\) of \(\mathrm{L}^p(\mathrm{G}, \beta)\) to be surjective ( \(1 \leqslant p < +\infty\) ), it is necessary and sufficient that there exist a \(c_p > 0\) such that \(\| \pmb{\mu}^{\vee} * g \|_q \geqslant c_p \| g \|_q\) for all \(g \in \mathrm{L}^q(\mathrm{G}, \beta)\) (where \(\frac{1}{p} + \frac{1}{q} = 1\) ) (cf.~TVS, IV, §4, No.~2, remarks following Cor. 3 of Th. 1).
\item
  If \(\mu\) is a measure with base \(\beta\) , and G is not discrete, show that \(\gamma(\mu)\) is never surjective for \(1 \leqslant p \leqslant +\infty\) (for \(p < +\infty\) , use \(a\) ), reducing to the case that the density of \(\mu\) with respect to \(\beta\) belongs to \(\mathcal{K}(\mathrm{G})\) ; for \(p = +\infty\) , argue directly by observing that for \(f \in \mathrm{L}^1(\mathrm{G}, \beta)\) and \(g \in \mathrm{L}^\infty(\mathrm{G}, \beta)\) , \(f * g\) is uniformly continuous for the right uniform structure).
\item
  If G is abelian, \(1 \leqslant p \leqslant 2\) and the endomorphism \(\pmb{\gamma}(\mu)\) of \(\mathrm{L}^p(\mathrm{G},\beta)\) is surjective, show that it is bijective (using regularization, show that if \(\pmb{\gamma}(\mu)\) is not injective then the endomorphism \(\pmb{\gamma}(\check{\mu})\) of \(\mathrm{L}^q(\mathrm{G},\beta)\) is not injective: one makes use of Ch. IV, §6, No.~5, Cor. of Prop. 4).
\item
  Take \(\mathbf{G} = \mathbf{Z}\) and \(\mu = \varepsilon_1 - \varepsilon_0\) ; show that \(\pmb{\gamma}(\mu)\) is injective in the \(\mathrm{L}^p (\mathbf{Z},\beta)\) for which \(p\neq +\infty\) , but not in \(\mathrm{L}^{\infty}(\mathbf{Z},\beta)\) , and it is not surjective for any value of \(p\) .
\end{enumerate}

¶ 18) Notations and conventions are those of Exer. 15.

\begin{enumerate}
\def\labelenumi{\alph{enumi})}
\item
  Let \(\mu\) be a bounded measure on \(G\) , whose support contains at least two distinct points. Show that there exists a compact set \(K\) , and a number \(k\) such that \(0 < k < 1\) , for which \(\| \varphi_{sK} \cdot \mu \| \leqslant k \| \mu \|\) for all \(s \in G\) . (If \(t, t'\) are two distinct points in the support of \(\mu\) , take for \(K\) a neighborhood of \(e\) sufficiently small that \(sK\) cannot intersect both \(tK\) and \(t'K\) .)
\item
  Let \(\mu\) be a bounded measure \(\geqslant 0\) on \(G\) , whose support contains at least two distinct points. Show that there exists a function \(f \in L^{\infty}(G, \beta)\) , not equivalent to a function \(\geqslant 0\) , such that \(\mu * f\) is \(\geqslant 0\) locally almost everywhere for \(\beta\) . (Use \(a\) ), taking \(f\) to be equal to \(-1\) on \(K\) and to a suitable positive constant on \(\mathbb{C}K\) .) If, in addition, the support of \(\mu\) is compact, there exists a function \(f\) having the preceding properties and having compact support.
\item
  Let \(\mu_{1},\mu_{2}\) be two nonzero, bounded positive measures on G, permutable for convolution, such that the support K of \(\mu_{1}\) is compact and contains \(e\) . Set \(\mu = \mu_{1} + \mu_{2}\) . Let V be a symmetric compact neighborhood of \(e\) containing K, \(g\) the function equal to \(\mu * \varphi_{\mathrm{V}}\) on \(\mathbb{C}\mathrm{V}\) , and to 0 on V; show that the function \(\mu * (g - (\mu_1 * \varphi_{\mathrm{V}}))\) is \(\geqslant 0\) locally almost everywhere in \(\mathbb{C}(\mathrm{V}^2)\) . On the other hand, show that there exists a compact set H such that the function \(\mu_{2} * \varphi_{\mathrm{H}}\) is \(\geqslant 0\) almost everywhere in \(\mathrm{V}^2\) .
\item
  Deduce from c) that if \(\mu\) is a bounded positive measure on G whose support contains at least two distinct points, there exists a function f that belongs to every \(L^{p}(G,\beta)\) ( \(1 \leqslant p \leqslant +\infty\) ), is not equivalent to a function \(\geqslant 0\) , and is such that \(\mu * f\) is \(\geqslant 0\) locally almost everywhere, when one makes in addition one of the following hypotheses: \(\alpha\) ) G is abelian; \(\beta\) ) there is a point \(a \in G\) such that \(\mu(\{a\}) > 0\) . (Suitably decompose \(\mu\) as a sum of two permutable measures \(\geqslant 0\) .)
\end{enumerate}

¶ 19) Notations and conventions are those of Exer. 15.

\begin{enumerate}
\def\labelenumi{\alph{enumi})}
\item
  Show that if a positive bounded measure \(\mu\) on \(G\) is such that for every function \(f \in L^{p}(G, \beta)\) ( \(p\) given, \(1 \leqslant p < \infty\) ) the relation \(\mu * f \geqslant 0\) locally almost everywhere implies \(f \geqslant 0\) locally almost everywhere, then one can conclude that \(\mu\) is a point measure in each of the following cases: \(\alpha)\) \(G\) is compact; \(\beta)\) \(G\) is discrete; \(\gamma)\) \(G\) is abelian (use Exer. 18) (*). For \(p = +\infty\) , the same conclusion is valid without a supplementary hypothesis on \(G\) .
\item
  Let \(\mu\) be a positive bounded measure on \(G\) such that: \(1^{\circ} \pmb{\gamma}(\mu)\) is a surjective endomorphism of \(L^{1}(G, \beta)\) ; \(2^{\circ}\) for every function \(f \in L^{1}(G, \beta)\) , the relation \(\mu * f \geqslant 0\) locally almost everywhere implies \(f \geqslant 0\) locally almost everywhere. Show that \(\mu\) is then a point measure. (Observe that for every function \(g \in L^{\infty}(G, \beta)\) , the relation \(\stackrel{\vee}{\mu} * g \geqslant 0\) locally almost everywhere implies \(g \geqslant 0\) locally almost everywhere.)
\end{enumerate}

\begin{enumerate}
\def\labelenumi{\arabic{enumi})}
\setcounter{enumi}{19}
\tightlist
\item
  Let \(G\) be a locally compact group, \(\beta\) a left Haar measure on \(G\) .
\end{enumerate}

\begin{enumerate}
\def\labelenumi{\alph{enumi})}
\item
  Let \(f\) be a bounded numerical function on \(G\) , uniformly continuous for the left uniform structure. Show that for every bounded measure \(\mu\) on \(G\) , \(\mu * f\) is uniformly continuous for the left uniform structure; moreover, with the notations of No.~7, Prop. 19, \(f\) is the limit, for the topology of uniform convergence in \(G\) , of the functions \(f_V * f\) with respect to the section filter of \(\mathfrak{B}\) .
\item
  Take \(\mathbf{G} = \mathbf{T}\) ; give an example of a continuous function \(h\) on \(\mathbf{T}\) that is not of the form \(f*g\) , where \(f\) and \(g\) belong to \(\mathrm{L}^2 (\mathbf{T},\beta)\) . (Use Exers. 15 b) and 16 of Ch. IV, §6.)
\end{enumerate}

¶ 21) Let G be a locally compact group, β a left Haar measure on G; one canonically identifies L¹(G,β) with a subspace of M¹(G).

\begin{enumerate}
\def\labelenumi{\alph{enumi})}
\item
  With the notations of §3, Exer. 11, let A be a relatively compact subset of \(\mathcal{M}_{\mathrm{III}}\) , B a subset of \(\mathrm{L}^1 (\mathrm{G},\beta)\) , relatively compact in \(\mathcal{M}_{\mathrm{IV}}\) . Show that the restriction to \(\mathrm{A}\times \mathrm{B}\) of the mapping \((\mu ,\nu)\mapsto \mu *\nu\) of \(\mathcal{M}_{\mathrm{III}}\times \mathcal{M}_{\mathrm{IV}}\) into \(\mathcal{M}_{\mathrm{IV}}\) is continuous. (First prove that if A and B are compact, the image of \(\mathrm{A}\times \mathrm{B}\) under this mapping is compact; use Exer. 10 of §3, as well as the criterion \(\alpha\) ) of Ch. V, §5, Exer. 15. To show next that \((\mu ,\nu)\mapsto \mu *\nu\) is continuous, make use of Exer. 11 c) of §3.)
\item
  Let \(\mathcal{U}_s^\infty (\mathrm{G})\) be the set of bounded numerical functions on \(\mathbf{G}\) uniformly continuous for the left uniform structure. Denote by \(\mathcal{T}_{\mathrm{II}}\) the topology \(\sigma (\mathcal{M}^1 (\mathrm{G}),\mathcal{U}_s^\infty (\mathrm{G}))\) and by \(\mathcal{M}_{\mathrm{II}}\) the space \(\mathcal{M}^1 (\mathrm{G})\) equipped with \(\mathcal{T}_{\mathrm{II}}\) . Taking \(\mathbf{G} = \mathbf{R}\) , give an example of a sequence \((\mu_n)\) of measures on \(\mathbf{G}\) tending to 0 for \(\mathcal{T}_{\mathrm{II}}\) and a sequence \((f_n)\) of functions in \(\mathrm{L}^1 (\mathrm{G},\beta)\) , tending to 0 for \(\mathcal{T}_{\mathrm{IV}}\) (or, what comes to the same, for the topology \(\sigma (\mathrm{L}^1 (\mathrm{G},\beta),\mathrm{L}^\infty (\mathrm{G},\beta))\) ), such that the sequence \((\mu_n*f_n)\) does not tend to 0 for \(\mathcal{T}_{\mathrm{IV}}\) (take \(f_{n}(t) = \sin nt\) in the interval \([0,\pi ]\) , \(f_{n}(t) = 0\) elsewhere).
\item
  Let A be a relatively compact subset of \(\mathcal{M}_{\mathrm{II}}\) , B a subset of \(\mathrm{L}^1 (\mathrm{G},\beta)\) , relatively compact for the topology of the norm \(\| f\| _1\) . Show that the restriction to \(\mathbf{A}\times \mathbf{B}\) of the mapping \((\mu ,f)\mapsto \mu *\textit{f}\) of \(\mathcal{M}_{\mathrm{III}}\times \mathrm{L}^1 (\mathrm{G},\beta)\) into \(\mathrm{L}^1 (\mathrm{G},\beta)\) is continuous (L \(^1 (\mathrm{G},\beta)\) being equipped with its normed space topology). (Reduce to proving that for \(f\in \mathcal{K}(\mathrm{G})\) , the mapping \(\mu \mapsto \mu *\textit{f}\) of \(\mathcal{M}_{\mathrm{II}}\) into \(\mathrm{L}^1 (\mathrm{G},\beta)\) is continuous in A. Using the fact that for \(g\in \mathrm{L}^{\infty}(\mathrm{G},\beta)\) and \(f\in \mathcal{K}(\mathrm{G})\) , \(g*\textit{f}\) is uniformly continuous for the left uniform structure, first show that the mapping \(\mu \mapsto \mu *\textit{f}\) of \(\mathcal{M}_{\mathrm{II}}\) into \(\mathcal{M}_{\mathrm{IV}}\) is continuous. Next, using Exer. 20 a), reduce to proving that if C is a subset of \(\mathrm{L}^1 (\mathrm{G},\beta)\) compact for the topology \(\sigma (\mathrm{L}^{1}(\mathrm{G},\beta),\mathrm{L}^{\infty}(\mathrm{G},\beta))\) and \(h\in \mathcal{K}(\mathbf{G})\) , the mapping \(\mu \mapsto \mu *h\) of C into \(\mathrm{L}^1 (\mathrm{G},\beta)\) is continuous when \(\mathrm{L}^1 (\mathrm{G},\beta)\) is equipped with its normed space topology. By the same reasoning as in \(a\) ), show that for this it suffices to prove that the image of C under this mapping is compact for the normed space topology. To this end, use Šmulian's theorem, the fact that C is cramped (§3, Exer. 10) and Lebesgue's theorem.)
\item
  Let A be a relatively compact subset of \(\mathcal{M}_{\mathrm{II}}\) containing 0, B a relatively compact subset of \(\mathcal{M}_{\mathrm{I}}\) ; show that for every \(\nu_0 \in \mathrm{B}\) , the restriction to \(\mathrm{A} \times \mathrm{B}\) of the mapping \((\mu, \nu) \mapsto \mu * \nu\) of \(\mathcal{M}_{\mathrm{II}} \times \mathcal{M}_{\mathrm{I}}\) into \(\mathcal{M}_{\mathrm{II}}\) is continuous at the point \((0, \nu_0)\) . (Make use of Exer. 20 a), and Exer. 21 c).
\item
  Let A be a bounded subset of \(\mathcal{M}^1 (\mathrm{G})\) , \(f\) a function in \(\mathrm{L}^1 (\mathrm{G},\beta)\) ; show that the restriction to A of the mapping \(\mu \mapsto \mu *f\) of \(\mathcal{M}_{\mathrm{III}}\) into \(\mathcal{M}_{\mathrm{IV}}\) is continuous.
\item
  Let \(\mathcal{M}_{+}^{1}(G)\) be the set of bounded positive measures on G. Show that if a subset A of \(\mathcal{M}_{+}^{1}(G)\) is compact for \(T_{II}\) , it is also compact for \(T_{III}\) . (Reduce to proving that A is a cramped set. Argue by contradiction, using the fact that for \(f \in \mathcal{K}(G)\) , the image of A under the mapping \(\mu \mapsto \mu * f\) is a cramped set by virtue of c).)
\end{enumerate}

\(g)\) Show that for \(\mathbf{G} = \mathbf{R}\) , the topologies induced on \(\mathcal{M}_{+}^{1}(\mathbf{G})\) by \(\mathcal{T}_{\mathrm{II}}\) and \(\mathcal{T}_{\mathrm{III}}\) are distinct.

¶ 22) Notations are those of Exer. 21 of §4 and Exer. 11 of §3.

\begin{enumerate}
\def\labelenumi{\alph{enumi})}
\item
  Let \((\mu_n)\) be a sequence of bounded measures on \(G\) . Assume that for every function \(f \in L^1(G, \beta)\) , the sequence \((\mu_n * f)\) converges to 0 in \(\mathcal{M}_I\) . Show that the sequence of norms \((\|\mu_n\|)\) is bounded (use the Banach-Steinhaus theorem for the family of mappings \(f \mapsto \mu_n * f\) of \(L^1(G, \beta)\) into itself, as well as Exer. 15). Deduce from this that the sequence \((\mu_n)\) tends to 0 in \(\mathcal{M}_I\) (use Exer. 20).
\item
  Let \((\mu_n)\) be a sequence of bounded measures on \(G\) such that the sequence \((\mu_n * f)\) tends vaguely to 0 for every function \(f \in L^1(G, \beta)\) ; show that the sequence \((\mu_n)\) tends vaguely to 0 (for every compact subset \(K\) of \(G\) , show, arguing as in \(a\) ), that the sequence \((|\mu_n|(K))\) is bounded). Give an example (with \(G = Z\) ) where the sequence \((\|\mu_n\|)\) is not bounded.
\item
  Let \((\mu_n)\) be a sequence of bounded measures on \(G\) such that for every function \(f \in L^1(G, \beta)\) , the sequence \((\mu_n * f)\) converges to 0 in \(\mathcal{M}_{II}\) ; show that the sequence \((\mu_n)\) converges to 0 in \(\mathcal{M}_{II}\) (make use of \(a\) )).
\end{enumerate}

\begin{enumerate}
\def\labelenumi{\arabic{enumi})}
\setcounter{enumi}{22}
\item
  Let \(G\) be a locally compact group, \(\beta\) a left Haar measure on \(G\) , \(f\) and \(g\) two positive locally \(\beta\) -integrable functions. Assume that \(f\) and \(g\) are convolvable and that one of these functions is zero on the complement of a countable union of compact sets. Show that \(f * g\) is equal locally almost everywhere to a lower semi-continuous function (use Prop. 15 of No.~5).
\item
  Show that the inequalities of Props. 12 and 15 of No.~5 cannot be improved by inserting a constant factor c \textless{} 1 in their right members.
\item
  Let G be a locally compact group, \(\beta\) a left Haar measure on G, \(\mu\) a positive measure on G. If A is a \(\mu\) -integrable set and B is a Borel set in G, the function
\end{enumerate}

\[
u: s \mapsto \mu (\mathrm{A} \cap s \mathrm{B})
\]

is \(\beta\) -measurable in G; if, in addition, \(B^{-1}\) is \(\beta\) -integrable, then so is \(u\) , and

\[
\int \mu (\mathrm{A} \cap s \mathrm{B}) d \beta (s) = \mu (\mathrm{A}) \beta (\mathrm{B} ^ {- 1}).
\]

Give an example of a measure \(\mu\) such that \(s \mapsto \mu(A \cap sB)\) is not continuous; if \(\mu\) is a measure with base \(\beta\) , and if A is \(\mu\) -integrable and B is a Borel set, then the function \(s \mapsto \mu(A \cap sB)\) is continuous on G.

¶ 26) Let G be a locally compact group, β a left Haar measure on G. For a subset H of L\^{}p(G,β) (1 ≤ p \textless{} +∞) to be relatively compact (for the topology of convergence in mean of order p), it is necessary and sufficient that the following conditions be satisfied: 1° H is bounded in L\^{}p(G,β); 2° for every ε \textgreater{} 0, there exists a compact subset K of G such that \textbar\textbar fφ\_G-K\textbar\textbar\_p ≤ ε for every function f ∈ H; 3° for every ε \textgreater{} 0, there exists a neighborhood V of e in G such that \textbar\textbar γ(s)f - f\textbar\textbar\_p ≤ ε for all f ∈ H and s ∈ V. (To prove that the conditions are sufficient, observe that if g ∈ K(G) and if L is a compact subset of G, then the image, under the mapping f ↦ g * f, of the set φ\_L · H, is an equicontinuous subset of K(G).)

¶ 27) Let G be a locally compact group, β a left Haar measure on G. If G is not reduced to e, then L¹(G,β) is an algebra (for convolution) admitting divisors of zero other than 0. One may proceed as follows to form two nonzero elements f, g of L¹(G,β) such that f * g = 0:

\(1^{\circ}\) The case that \(\mathbf{G}\) contains a compact subgroup \(\mathbf{H}\) , not reduced to \(e\) , in which \(\Delta(x)=1\) . Take \(f\) to be the characteristic function of a set \(\mathbf{A}\) , and \(g\) to be a difference \(\varphi_{s\mathbf{B}}-\varphi_{\mathbf{B}}\) of two characteristic functions, with \(\mathbf{A}\) and \(\mathbf{B}\) suitably chosen.

\(2^{\circ}\) The case that \(\mathbf{G} = \mathbf{Z}\) . Show that one can then take

\[
f (n) = \frac {1}{2 n - 1} - \frac {1}{2 n + 1}
\]

for all \(n \in \mathbf{Z}\) , and \(g(n) = f(-n)\) .

\(3^{\circ}\) The general case. First prove that there exists an \(a \neq e\) in G such that \(\Delta(a) = 1\) . The closure H in G of the subgroup generated by \(a\) is then either a compact subgroup or a subgroup isomorphic to Z (GT, V, §1, Exer. 2). In the first case, use the result of \(1^{\circ}\) ; in the second, take

\[
f (t) = \sum_ {n = - \infty} ^ {+ \infty} \alpha_ {n} \varphi_ {\mathrm{U} a ^ {- n}} (t), \quad g (t) = \sum_ {n = - \infty} ^ {+ \infty} \beta_ {n} \varphi_ {a ^ {n} \mathrm{U}} (t),
\]

suitably choosing U and the sequences \((\alpha_{n})\) , \((\beta_{n})\) with the help of \(2^{\circ}\) .

\begin{enumerate}
\def\labelenumi{\arabic{enumi})}
\setcounter{enumi}{27}
\tightlist
\item
  Let \(G_1, G_2\) be two locally compact groups, \(\beta_1, \beta_2\) left Haar measures on \(G_1, G_2\) , and \(A_1\) (resp. \(A_2\) ) the topological algebra (over \(R\) ) \(L^1(G_1, \beta_1)\) (resp. \(L^1(G_2, \beta_2)\) ). Let \(T\) be an algebra isomorphism of \(A_1\) onto \(A_2\) , such that the relation \(f \geqslant 0\) almost everywhere is equivalent to \(T(f) \geqslant 0\) almost everywhere.
\end{enumerate}

\begin{enumerate}
\def\labelenumi{\alph{enumi})}
\item
  Show that \(\mathbf{T}\) is an isomorphism of topological algebras. (First note that if a decreasing sequence \((f_n)\) of elements of \(\mathbf{A}_1\) tends to 0 almost everywhere, then the same is true of the sequence of the \(\mathrm{T}(f_n)\) ; deduce from this that for every sequence \((f_n)\) that is bounded above in \(\mathbf{A}_1\) , \(\mathrm{T}(\sup (f_n)) = \sup (\mathrm{T}(f_n))\) , and conclude with the help of Fatou's lemma.)
\item
  Show that T may be extended, in only one way, to an isomorphism of the topological algebra \(\mathcal{M}^{1}(G_{1})\) onto the topological algebra \(\mathcal{M}^{1}(G_{2})\) , and that there exist a topological isomorphism u of \(G_{1}\) onto \(G_{2}\) and a continuous homomorphism \(\chi\) of \(G_{1}\) into \(R_{+}^{*}\) such that \(\mathrm{T}(\varepsilon_{s}) = \chi(s)\varepsilon_{u(s)}\) (make use of Exer. 15 a) and Exer. 19 b); to prove that u and \(\chi\) are continuous, observe that T defines an isomorphism of the algebra \(\mathcal{L}(A_{1})\) of continuous endomorphisms of the topological vector space \(A_{1}\) onto the analogous algebra \(\mathcal{L}(A_{2})\) , and that this isomorphism is bicontinuous for the topology of pointwise convergence; on the other hand, observe that \(s \mapsto \delta(s^{-1})\) is an isomorphism of \(G_{1}\) onto a multiplicative subgroup of \(\mathcal{L}(A_{1})\) .
\item
  Deduce from \(b)\) that, for every \(f \in \mathbf{A}_1\) ,
\end{enumerate}

\[
(T (f)) (t) = \chi (u ^ {- 1} (t)) f (u ^ {- 1} (t))
\]

for all \(t \in \mathbf{G}_2\) .

Recall (A, III, §2, Exer. 9) that there exist finite groups \(\mathbf{G}_1, \mathbf{G}_2\) such that the algebras \(\mathrm{L}^1(\mathrm{G}_1, \beta_1)\) and \(\mathrm{L}^1(\mathrm{G}_2, \beta_2)\) are isomorphic without \(\mathbf{G}_1\) and \(\mathbf{G}_2\) being so.

